%-*-latex-*-
%-*-latex-*-
\newcommand\COURSE{ciss350}
\newcommand\ASSESSMENT{q10706}
\newcommand\ASSESSMENTTYPE{Quiz}
\newcommand\POINTS{\textwhite{xxx/xxx}}

\input{myquizpreamble}
\input{yliow}
\input{\COURSE}
%-*-latex-*-

%-*-latex-*-

%-*-latex-*-
\renewcommand\TITLE{\ASSESSMENTTYPE \ \ASSESSMENT}

\renewcommand\EMAIL{}
\renewcommand\AUTHOR{eagraham1@cougars.ccis.edu}

\textwidth=6in
\begin{document}
\topmatterthree

For the following questions,
here's an example on how to describe a binary tree in text.
For this tree
\begin{console}[fontsize=\footnotesize]
            (10)
          /      \
      ( 4)        (25)
     /           /    \
 ( 1)        (12)      (27)
                           \
                            (29)
\end{console}
you write
\begin{console}[fontsize=\footnotesize]
[10, [4, 25]]
[4, [1, None]]
[25, [12, 27]]
[1, [None, None]]
[12, [None, None]]
[27, [None, 29]]
[29, [None, None]]
\end{console}

%------------------------------------------------------------------------------
\nextq
For this tree:
\begin{console}[fontsize=\footnotesize]
            (10)
          /      \
      ( 4)        (25)
     /           /    \
 ( 1)        (12)      (27)
                           \
                            (29)
\end{console}
What is the tree after performing a left rotation at the node
with key value of $10$?
\\
\ANSWER
\begin{answercode}
[25, [10, 27]]
[10, [4, 12]]
[27, [None, 29]]
[4, [1, None]]
[12, [None, None]]
[29, [None, None]]
[1, [None, None]]
\end{answercode}

%------------------------------------------------------------------------------
\nextq
For this tree:
\begin{console}[fontsize=\footnotesize]
            (10)
          /      \
      ( 4)        (25)
     /           /    \
 ( 1)        (12)      (27)
                           \
                            (29)
\end{console}
What is the tree after performing a right rotation at the node
with key value of $10$?
\\
\ANSWER
\begin{answercode}
[4, [1, 10]]
[1, [None, None]]
[10, [None, 25]]
[25, [12, 27]]
[12, [None, None]]
[27, [None, 29]]
[29, [None, None]]
\end{answercode}

%------------------------------------------------------------------------------
\nextq
For this BST:
\begin{console}[fontsize=\footnotesize]
            (10)
          /      \
      ( 4)        (25)
     /           /    \
 ( 1)        (12)      (27)
                           \
                            (29)
\end{console}
insert 30 and then balance it (if necessary).
Write the resulting tree in text below.
\ANSWER
\begin{answercode}
[25, [10, 29]]
[10, [4, 12]]
[29, [27, 30]]
[4, [1, None]]
[12, [None, None]]
[27, [None, None]]
[30, [None, None]]
[1, [None, None]]
\end{answercode}

%------------------------------------------------------------------------------
\nextq
For this BST:
\begin{console}[fontsize=\footnotesize]
            (10)
          /      \
      ( 4)        (25)
     /           /    \
 ( 1)        (12)      (27)
                           \
                            (29)
\end{console}
insert 2 and then balance it (if necessary).
Write the resulting tree in text below.
\ANSWER
\begin{answercode}
[10, [4, 25]]
[2, [1, 4]]
[25, [12, 27]]
[1, [None, None]]
[4, [None, None]]
[12, [None, None]]
[27, [None, 29]]
[29, [None, None]]
\end{answercode}

%------------------------------------------------------------------------------
%-*-latex-*-
\newpage
\input{instructions}
\end{document}

